% !TEX program = xelatex
% Deck (cards format) — hand-authored storybook deck of the laer-klokken week of
% 14 September 2026. Sibling of ../laer-klokken-sept-2026/ (the broadside);
% same facts, addressed to the klokkers, one scene per card.
\documentclass[11pt]{article}

\usepackage{fontspec}
\usepackage{unicode-math}
% Latin Modern by OTF filename (not font name): xelatex on macOS cannot
% resolve the family by name and silently produces a broken card PDF.
\setmainfont{lmroman10-regular.otf}[
  BoldFont=lmroman10-bold.otf,
  ItalicFont=lmroman10-italic.otf,
  BoldItalicFont=lmroman10-bolditalic.otf,
]
\setsansfont{lmsans10-regular.otf}[
  BoldFont=lmsans10-bold.otf,
  ItalicFont=lmsans10-oblique.otf,
  BoldItalicFont=lmsans10-boldoblique.otf,
]
\setmonofont{lmmono10-regular.otf}[Scale=0.88,
  BoldFont=lmmonolt10-bold.otf,
  ItalicFont=lmmono10-italic.otf,
]

\usepackage{graphicx}
\graphicspath{{figures/}{figures/illys/}{figures/paintings/}{figures/klokkers/}}
\usepackage{booktabs}
\usepackage{tabularx}
\usepackage{adjustbox}
\usepackage{ragged2e}
\usepackage{microtype}

\makeatletter
\def\input@path{{../}}
\makeatother
\usepackage{ac-paper-cards}
% the cards style sets \ac in small caps with a plain period; the brand wants the pink dot.
\renewcommand{\ac}{Aesthetic{\acdot}Computer}
\newcommand{\dk}[1]{{\itshape #1}}
\newcommand{\illy}[1]{\begin{center}\includegraphics[width=0.5\textwidth]{#1}\end{center}\vspace{-0.4em}}
% chat syntax, carried into print: @handle is a tappable profile link, #code opens the painting.
% chat paints handles pink ([255,150,200] on its dark ground); this is that pink at paper contrast.
\definecolor{achandle}{RGB}{206,72,132}
\newcommand{\h}[1]{\href{https://prompt.ac/@#1}{\color{achandle}\textbf{@#1}}}
\newcommand{\hh}[2]{\href{https://prompt.ac/@#1}{\color{achandle}\textbf{@#2}}}
\newcommand{\pc}[1]{\href{https://aesthetic.computer/media/paintings/#1.png}{\color{acpurple}\texttt{\##1}}}
% a painting tile: image, then its code and painter in tiny type
\newcommand{\tile}[3]{\begin{minipage}[t]{0.29\textwidth}\centering\includegraphics[width=\linewidth]{#1}\\[-0.1em]{\tiny\pc{#2}\ \h{#3}}\end{minipage}}
\newcommand{\hhtile}{\begin{minipage}[t]{0.184\textwidth}\centering
\includegraphics[width=\linewidth]{FM_Vigil}\\[-0.15em]{\fontsize{5pt}{6pt}\selectfont\hh{FM_Vigil}{FM\_Vigil}}\end{minipage}}
\newcommand{\ktile}[2]{\begin{minipage}[t]{0.184\textwidth}\centering\includegraphics[width=\linewidth]{#1}\\[-0.15em]{\fontsize{5pt}{6pt}\selectfont\h{#2}}\end{minipage}}

\hypersetup{
  pdftitle={Laer Klokken Escapades --- the week of 14 September 2026, on cards},
}

\renewcommand{\acpdfbase}{laer-klokken-escapades-26}
\acdraftfalse
\begin{document}

% ============================================================
% TITLE CARD
% ============================================================
\thispagestyle{empty}
\vspace*{0.05em}
\begin{center}
\href{https://laklok.com}{\includegraphics[width=0.66\textwidth]{cover}}\par\vspace{0.25em}
{\acbold\fontsize{18pt}{22pt}\selectfont\color{acdark} Laer Klokken Escapades}\par
\vspace{0.1em}
{\fontsize{9pt}{11pt}\selectfont\color{acpink} the week of 14 september 2026, on cards --- for the klokkers, before stubbekøbing}\par
\vspace{0.2em}
{\scriptsize\color{acgray}\h{jeffrey}\,\textperiodcentered\,\ac{}\,\textperiodcentered\,laklok.com}\par
\vspace{0.2em}
\rule{0.5\textwidth}{0.5pt}\par
\vspace{0.12em}
{\scriptsize\color{acgray} tuesday 15 september 2026 \,\textperiodcentered\, read off chat-clock and the git log}\par
\end{center}

% ============================================================
% INDEX CARD
% ============================================================
\cardindex

% ============================================================
% LETTER CARD
% ============================================================
\clearpage
\begin{accentcard}
\cardtitle{dear klokkers}

\noindent here is your week, as i read it off the chat on tuesday --- where you're going on saturday, what happened in the clock channel, and what \ac{} handed you along the way. print it, cut it, pin it up, hand it round at stubbekøbing. the last card is a question. \hfill\textit{---\,\h{jeffrey}}
\end{accentcard}

% ============================================================
% BODY
% ============================================================
\section{the trip}

\illy{trip}

folkets folkemøde, stubbekøbing, falster --- a people's meeting without lobbyists, one saturday a year. this year it has a clock in it.

\textbf{saturday 19 september, 13:30 to 14:30, rådhuset.} \dk{laer klokken med goodiepal} is in the printed program. their blurb: the composer and clockmaker records a new episode of his unpredictable radio show, about quitting social media for good and building your own computing on \ac{}. visitors are told to try the commands and catch up on old episodes. our pitch, written by someone else --- the best kind.

\section{the klokkers}

the newest painting on each of your profiles, tuesday afternoon. tap a name to visit.

\begin{center}
\ktile{prutti}{prutti}\hfill\ktile{HendeFraHoeve}{HendeFraHoeve}\hfill\ktile{snakes}{snakes}\hfill\ktile{YouthOrgan}{YouthOrgan}\hfill\ktile{tyttebaer}{tyttebaer}\\[0.9em]
\ktile{osamakimlarsen}{osamakimlarsen}\hfill\ktile{q13R608tx}{q13R608tx}\hfill\ktile{chigichigi}{chigichigi}\hfill\hhtile\hfill\ktile{alki}{alki}
\end{center}

\h{oliverstilling} and \h{Wallowing} have no paintings yet. that is a hint, not a complaint.

\section{the plan, seven points}

\illy{plan}

posted by \h{prutti} on tuesday at 09:08, numbering his own:

\begin{enumerate}
\item record a laer klokken program live.
\item every klokker in the room gets a segment.
\item the show runs 56 minutes.
\item there is no money, because there just isn't any.
\item a reception on the HVAD record.
\item a bed in stubbekøbing the night before and after.
\item a bus down from copenhagen --- if you say you're coming.
\end{enumerate}

coda: \dk{det bliver super duper}. it will.

\section{the phone line}

\illy{phone}

\dk{ring til Torras telefon. nu.} then \dk{linjen er åben nu} at 10:58, and again at 11:39. \h{tyttebaer} and \h{YouthOrgan} called in while the rest of the channel watched. the clock program is being made live, over a phone, in public, on a tuesday morning. that's the format.

\textbf{who's in so far:} \h{HendeFraHoeve} first (\dk{jeg skal med!}), then \h{snakes}, \h{YouthOrgan}, \h{tyttebaer} --- and \h{oliverstilling}, who said back in august he'd see goodie there. \h{alki} wants \dk{bus intel}. \h{osamakimlarsen} wants to meet up beforehand \dk{og se hvem er folket}. prutti's reply to each: \dk{tjek}.

\section{the chat became a video}

\illy{video}

\hh{FM_Vigil}{FM\_Vigil} recorded the chat and put it on youtube. \h{HendeFraHoeve} cut it three times in one afternoon, chasing bigger text --- \dk{der kommer straks en version med større chat}, then \dk{større men mere uklar}, then \dk{alle gode gange 3}.

prutti watched on an ipad at glasmagasinet in oslo: \dk{chatten er lidt lille men den er fucking fed}.

the room is editing itself now. i just have to keep it running --- and yes, a bigger-chat tema is coming, so the third re-render isn't needed next time.

\section{codes on the wire}

the paintings you dropped bare into the channel this month, as codes --- three characters, no preview, tap it and see. that was the whole idea.

\begin{center}
\tile{swk}{swk}{q13R608tx}\hfill\tile{w5o}{w5o}{osamakimlarsen}\hfill\tile{n5j}{n5j}{osamakimlarsen}\\[0.5em]
\tile{v5x}{v5x}{osamakimlarsen}\hfill\tile{ie6}{ie6}{osamakimlarsen}\hfill\tile{6x3}{6x3}{osamakimlarsen}\\[0.5em]
\tile{aue}{aue}{q13R608tx}\hfill\tile{jzg}{jzg}{q13R608tx}\hfill\tile{4ne}{4ne}{q13R608tx}
\end{center}

\section{datamusik provenance}

\illy{solder}

\h{q13R608tx} and \h{snakes} spent monday evening on who really made the datamusik records and the V/Vm audio --- ``jansky noise has also confirmed that datamusik 13 and 03 are both made by you.'' snakes: he made more sounds than he remembers.

q13 is \dk{soldering a bunch of wires once more to create a product}, and thanked osama for the stickers and the collage. the archive argues with itself. good.

\section{what the platform handed you}

everything here landed in the last thirty days, most of it built for this room.

\textbf{two laklok sisters, one census} (09-13). the raster laklok piece and the vector laklok.com/html now stay in lockstep, and both report which tema each visitor wears. a \texttt{realtime} tema drifts with the calendar.

\textbf{prutti's words learned to speak} (09-08). a vox chip on every \h{prutti} message --- his consented elevenlabs voice reads the line back, word by word. since 09-13 any signed-in handle can send a text through \texttt{say}.

\textbf{laer-klokken TV, and the bridge back} (09-01). a 24/7 youtube stream of the clock room; youtube viewers appear in chat-clock as televised guests.

\textbf{two things to check before saturday.} \h{snakes} said on 09-04 that goodiepal could not log in --- the host of the show should be able to get into his own show. and \h{kitten.cat} hit a ``verify'' prompt that didn't explain itself.

\section{every klokker is an email address}

\illy{mail}

\texttt{prutti@aesthetic.computer} exists. so does yours: your handle, at aesthetic.computer. the door is open, SPF and DKIM are on, and nobody had told the room until now.

prutti asked on 09-08 where to send three files ``for computer change'' and again on 09-11 --- \dk{what email adress should i use to file things to you?} letters from outside so far: zero. that address is the answer.

\h{q13R608tx} worked it out on his own: ``mail is the new private messaging system basically... but even better because its mail!!'' he's right.

\section{the road to falster}

\illy{road}

\begin{itemize}
\item \textbf{mon 14} --- the chat is recorded and cut into a video, three times.
\item \textbf{tue 15} --- the plan is posted and the phone line opens.
\item \textbf{fri 18} --- prutti comes down from norway to copenhagen. the festival sells bus tickets from there; prutti's own bus needs a headcount.
\item \textbf{sat 19, 13:30} --- rådhuset, stubbekøbing: 56 minutes recorded live.
\item \textbf{after} --- the episode joins the archive online.
\end{itemize}

\section{for saturday}

\illy{bell}

what i can bring to the room without being in it: this deck, printed; a QR to laklok.com on the door; the TV stream running on any screen in the hall; a mail address for every klokker who signs up on the day; a bigger-chat tema for whoever films it.

what else do you need for the gig? a card, a poster, a sound, a piece that only runs in stubbekøbing --- say it in chat-clock and it goes on the list.

\vspace{0.5em}
\noindent{\small\color{acgray}\dk{vi ses i laer-klokken.}}

\end{document}
